\chapter{Glyph curation}\label{module:curate}
\Q{Evidence}{Current implementation boundary}
\path{Software/DemoCode/curate.py}. Exact file hashes are in the accompanying source manifest. Names, declarations and source locations below come from the captured source, not a runtime import.
\Q{Purpose}{Input-to-output reading flow}
Inspect the stored profile and its variants, choose or remove entries, then persist the curated profile. This is profile preparation rather than inference-time model training.


\Q{Evidence}{Selected current module settings}
These are declarations or expressions, not recovered environment values. Source comments do not establish a measured optimum. Larger glyph coordinate tables are designed data in the bundled source.
\begin{lstlisting}[style=plain]
HERE = Path(__file__).resolve().parent
OUT = HERE / 'results' / 'glyph_curation'
\end{lstlisting}
\section*{Stage: Profile curation}
\section{Method background: Profile curation}\label{method:curate}
\Q{Purpose}{The role in this system}
Remove selected bad glyphs and rebuild representatives without training a new classifier.


\Q{Justification}{The reason for this construction}
Explicit variant indices and a saved curated profile make manual quality selection inspectable. A classifier label remains the base writer, not a new class for the suffix.


\Q{Alternatives}{What changes if another method is used}
Automatic prior/consensus filtering scales better but can retain systematic bad cuts. Human selection adds judgement and should not be hidden when interpreting results.


\Q{Library}{Implementation speed and accuracy}
Pillow draws inspection sheets; geometry and pool policy remain project-specific. A library UI or clustering tool changes workflow rather than proving better style. \cite{code2026,haines2016,pillowdraw}.


\section{ParseSpec}\label{fn:curate:ParseSpec}
\begin{lstlisting}[style=sig]
def ParseSpec(spec)
\end{lstlisting}
\Q{Evidence}{Source and method link}
\path{Software/DemoCode/curate.py} line 42. Implementation evidence: \cite{code2026}. Related method, alternatives and library discussion: Section~\ref{method:curate}.


\Q{Purpose}{What this contributes}
Remove selected bad glyphs and rebuild representatives without training a new classifier.


\Q{Method}{Current implementation}
Remove selected bad glyphs and rebuild representatives without training a new classifier. The excerpt below is from the current definition; the complete function is in the source snapshot. Its constants and branches, rather than an older signature, govern this path.


\begin{lstlisting}[style=plain]
def ParseSpec(spec):
    out = {}
    for part in (spec or "").split():
        if ":" not in part:
            raise SystemExit("expected char:indices, got %r" % part)
        ch, idx = part.split(":", 1)
        if len(ch) != 1:
            raise SystemExit("one character per entry, got %r" % ch)
        out.setdefault(ch, set()).update(int(i) for i in idx.split(",") if i != "")
    return out
\end{lstlisting}
\Q{Justification}{Interpret the method rather than infer a benchmark}
Explicit variant indices and a saved curated profile make manual quality selection inspectable. A classifier label remains the base writer, not a new class for the suffix. \cite{code2026,haines2016,pillowdraw}.


\section{LoadProfile}\label{fn:curate:LoadProfile}
\begin{lstlisting}[style=sig]
def LoadProfile(writer, curated=False)
\end{lstlisting}
\Q{Evidence}{Source and method link}
\path{Software/DemoCode/curate.py} line 55. Implementation evidence: \cite{code2026}. Related method, alternatives and library discussion: Section~\ref{method:curate}.


\Q{Purpose}{What this contributes}
Remove selected bad glyphs and rebuild representatives without training a new classifier.


\Q{Method}{Current implementation}
Remove selected bad glyphs and rebuild representatives without training a new classifier. The excerpt below is from the current definition; the complete function is in the source snapshot. Its constants and branches, rather than an older signature, govern this path.


\begin{lstlisting}[style=plain]
def LoadProfile(writer, curated=False):
    name = "%s_curated.json" % writer if curated else "%s.json" % writer
    path = paths.Profiles/name
    if not path.exists():
        raise SystemExit("no profile at %s" % path)
    return json.load(open(path, encoding="utf-8")), path
\end{lstlisting}
\Q{Justification}{Interpret the method rather than infer a benchmark}
Explicit variant indices and a saved curated profile make manual quality selection inspectable. A classifier label remains the base writer, not a new class for the suffix. \cite{code2026,haines2016,pillowdraw}.


\section{Listing}\label{fn:curate:Listing}
\begin{lstlisting}[style=sig]
def Listing(writer, curated=False)
\end{lstlisting}
\Q{Evidence}{Source and method link}
\path{Software/DemoCode/curate.py} line 132. Implementation evidence: \cite{code2026}. Related method, alternatives and library discussion: Section~\ref{method:curate}.


\Q{Purpose}{What this contributes}
Remove selected bad glyphs and rebuild representatives without training a new classifier.


\Q{Method}{Current implementation}
Remove selected bad glyphs and rebuild representatives without training a new classifier. The excerpt below is from the current definition; the complete function is in the source snapshot. Its constants and branches, rather than an older signature, govern this path.


\begin{lstlisting}[style=plain]
def Listing(writer, curated=False):
    prof, path = LoadProfile(writer, curated)
    glyphs = prof.get("glyphs") or {}
    print("%s, from %s: %d characters, %d variants in total"
          % (writer, path.name, len(glyphs), sum(len(v) for v in glyphs.values())))
    for ch in sorted(glyphs):
        vs = glyphs[ch]
        pd = [float(g.get("priorD", 0.5)) for g in vs]
        proto = sum(1 for g in vs if g.get("denoised"))
        print("  %-4r %3d variants   prior distance min %.2f median %.2f max %.2f%s"
              % (ch, len(vs), min(pd), float(np.median(pd)), max(pd),
                 "   (%d prototype)" % proto if proto else ""))
\end{lstlisting}
\Q{Justification}{Interpret the method rather than infer a benchmark}
Explicit variant indices and a saved curated profile make manual quality selection inspectable. A classifier label remains the base writer, not a new class for the suffix. \cite{code2026,haines2016,pillowdraw}.


\section{Variants}\label{fn:curate:Variants}
\begin{lstlisting}[style=sig]
def Variants(writer, curated=False)
\end{lstlisting}
\Q{Evidence}{Source and method link}
\path{Software/DemoCode/curate.py} line 198. Implementation evidence: \cite{code2026}. Related method, alternatives and library discussion: Section~\ref{method:curate}.


\Q{Purpose}{What this contributes}
The whole set in a form something else can draw the command line draws its own sheets, but the curation page in the web app needs the strokes themselves so it can draw them and be clicked on. the indices here are the ones Apply takes, so what is seen and what is removed cannot drift apart.


\Q{Method}{Current implementation}
Remove selected bad glyphs and rebuild representatives without training a new classifier. The excerpt below is from the current definition; the complete function is in the source snapshot. Its constants and branches, rather than an older signature, govern this path.


\begin{lstlisting}[style=plain]
def Variants(writer, curated=False):
    prof, path = LoadProfile(writer, curated)
    glyphs = prof.get("glyphs") or {}
    chars = []
    for ch in sorted(glyphs):
        vs = []
        for i, g in enumerate(glyphs[ch]):
            vs.append(dict(index=i,
                           priorD=round(float(g.get("priorD", 0.5)), 4),
                           advance=round(float(g.get("advance", 0.9)), 3),
                           denoised=bool(g.get("denoised")),
                           strokes=[[[round(float(p[0]), 4), round(float(p[1]), 4)]
                                     for p in s] for s in (g.get("strokes") or [])]))
        chars.append(dict(ch=ch, variants=vs))
    return dict(writer=writer, source=path.name, characters=chars,
                nChars=len(chars), nVariants=sum(len(c["variants"]) for c in chars),
                curatedFrom=prof.get("curatedFrom"),
                hasCurated=(paths.Profiles/("%s_curated.json" % writer)).exists())
\end{lstlisting}
\Q{Justification}{Interpret the method rather than infer a benchmark}
Explicit variant indices and a saved curated profile make manual quality selection inspectable. A classifier label remains the base writer, not a new class for the suffix. \cite{code2026,haines2016,pillowdraw}.


\section{Drop}\label{fn:curate:Drop}
\begin{lstlisting}[style=sig]
def Drop(writer, dropSpec=None, keepSpec=None, curated=False, outName=None)
\end{lstlisting}
\Q{Evidence}{Source and method link}
\path{Software/DemoCode/curate.py} line 219. Implementation evidence: \cite{code2026}. Related method, alternatives and library discussion: Section~\ref{method:curate}.


\Q{Purpose}{What this contributes}
Remove selected bad glyphs and rebuild representatives without training a new classifier.


\Q{Method}{Current implementation}
Remove selected bad glyphs and rebuild representatives without training a new classifier. The excerpt below is from the current definition; the complete function is in the source snapshot. Its constants and branches, rather than an older signature, govern this path.


\begin{lstlisting}[style=plain]
def Drop(writer, dropSpec=None, keepSpec=None, curated=False, outName=None):
    return Apply(writer, ParseSpec(dropSpec), ParseSpec(keepSpec), curated, outName)
\end{lstlisting}
\Q{Justification}{Interpret the method rather than infer a benchmark}
Explicit variant indices and a saved curated profile make manual quality selection inspectable. A classifier label remains the base writer, not a new class for the suffix. \cite{code2026,haines2016,pillowdraw}.


\section{Apply}\label{fn:curate:Apply}
\begin{lstlisting}[style=sig]
def Apply(writer, drop=None, keep=None, curated=False, outName=None)
\end{lstlisting}
\Q{Evidence}{Source and method link}
\path{Software/DemoCode/curate.py} line 224. Implementation evidence: \cite{code2026}. Related method, alternatives and library discussion: Section~\ref{method:curate}.


\Q{Purpose}{What this contributes}
Remove selected bad glyphs and rebuild representatives without training a new classifier.


\Q{Method}{Current implementation}
Remove selected bad glyphs and rebuild representatives without training a new classifier. The excerpt below is from the current definition; the complete function is in the source snapshot. Its constants and branches, rather than an older signature, govern this path.


\begin{lstlisting}[style=plain]
def Apply(writer, drop=None, keep=None, curated=False, outName=None):
    prof, path = LoadProfile(writer, curated)
    glyphs = dict(prof.get("glyphs") or {})
    drop = {k: set(v) for k, v in (drop or {}).items() if v}
    keep = {k: set(v) for k, v in (keep or {}).items()}
    for ch in list(drop)+list(keep):
        if ch not in glyphs:
            raise SystemExit("%r has no variants in %s" % (ch, path.name))

    removed, emptied = 0, []
    for ch, vs in list(glyphs.items()):
        if ch in keep:
            chosen = [g for i, g in enumerate(vs) if i in keep[ch]]
        elif ch in drop:
            chosen = [g for i, g in enumerate(vs) if i not in drop[ch]]
        else:
            continue
        removed += len(vs)-len(chosen)
        if not chosen:
            emptied.append(ch)
            del glyphs[ch]
        else:
            glyphs[ch] = chosen

    prof["glyphs"] = glyphs
    ideal = dict(prof.get("idealGlyphs") or {})
    medoid = dict(prof.get("idealMedoid") or {})
    adv = dict(prof.get("letterAdvance") or {})
    touched = set(drop) | set(keep)
    for ch in touched:
        if ch not in glyphs:
            ideal.pop(ch, None)
            medoid.pop(ch, None)
            adv.pop(ch, None)
            continue
        vs = glyphs[ch]
        p = Prototype(ch, vs)
        m = Medoid(ch, vs)
        if p is None:
            ideal.pop(ch, None)
        else:
            ideal[ch] = p
        if m is None:
            medoid.pop(ch, None)
        else:
            medoid[ch] = m
        adv[ch] = round(float(np.median([g["advance"] for g in vs])), 3)
    prof["idealGlyphs"] = ideal
    prof["idealMedoid"] = medoid
    prof["letterAdvance"] = adv
    prof["curatedFrom"] = path.name
    prof["curatedRemoved"] = removed
    prof["curatedNote"] = ("glyphs, idealGlyphs, idealMedoid and letterAdvance "
                           "were rebuilt from the kept variants; connectedness "
                           "is carried over unchanged because the adjacency it "
                           "was counted from is not stored")

    out = paths.Profiles/(outName or ("%s_curated.json" % writer))
# Excerpt ends; full definition continues in bundled source.
\end{lstlisting}
\Q{Justification}{Interpret the method rather than infer a benchmark}
Explicit variant indices and a saved curated profile make manual quality selection inspectable. A classifier label remains the base writer, not a new class for the suffix. \cite{code2026,haines2016,pillowdraw}.


\section{Main}\label{fn:curate:Main}
\begin{lstlisting}[style=sig]
def Main()
\end{lstlisting}
\Q{Evidence}{Source and method link}
\path{Software/DemoCode/curate.py} line 303. Implementation evidence: \cite{code2026}. Related method, alternatives and library discussion: Section~\ref{method:curate}.


\Q{Purpose}{What this contributes}
Remove selected bad glyphs and rebuild representatives without training a new classifier.


\Q{Method}{Current implementation}
Remove selected bad glyphs and rebuild representatives without training a new classifier. The excerpt below is from the current definition; the complete function is in the source snapshot. Its constants and branches, rather than an older signature, govern this path.


\begin{lstlisting}[style=plain]
def Main():
    ap = argparse.ArgumentParser(description=__doc__.splitlines()[0])
    ap.add_argument("action", choices=["sheet", "list", "drop"])
    ap.add_argument("writer")
    ap.add_argument("--drop", default=None, help='e.g. "a:2,5 e:0"')
    ap.add_argument("--keep", default=None, help='e.g. "a:0,1" keeps only those')
    ap.add_argument("--curated", action="store_true",
                    help="work from <writer>_curated.json instead of the original")
    ap.add_argument("--out", default=None, help="name for the written profile")
    a = ap.parse_args()
    if a.action == "sheet":
        Sheet(a.writer, a.curated)
    elif a.action == "list":
        Listing(a.writer, a.curated)
    else:
        if not (a.drop or a.keep):
            raise SystemExit("give --drop or --keep")
        Drop(a.writer, a.drop, a.keep, a.curated, a.out)
\end{lstlisting}
\Q{Justification}{Interpret the method rather than infer a benchmark}
Explicit variant indices and a saved curated profile make manual quality selection inspectable. A classifier label remains the base writer, not a new class for the suffix. \cite{code2026,haines2016,pillowdraw}.


\section*{Stage: Consensus medoids and prototypes}
Method rationale, alternatives and library comparison: Section~\ref{method:medoid}.
\section{Medoid}\label{fn:curate:Medoid}
\begin{lstlisting}[style=sig]
def Medoid(ch, vs)
\end{lstlisting}
\Q{Evidence}{Source and method link}
\path{Software/DemoCode/curate.py} line 153. Implementation evidence: \cite{code2026}. Related method, alternatives and library discussion: Section~\ref{method:medoid}.


\Q{Purpose}{What this contributes}
REBUILDING WHAT THE SET DECIDES the library is not the only thing that comes from the variants: the prototype letterform, the most typical real variant and the letter advances are all derived from them, so all three are recomputed rather than left describing glyphs that are no longer there.


\Q{Method}{Current implementation}
Choose representative writer glyphs while rejecting fragmented outliers. The excerpt below is from the current definition; the complete function is in the source snapshot. Its constants and branches, rather than an older signature, govern this path.


\begin{lstlisting}[style=plain]
def Medoid(ch, vs):
    grids = [(ST._ShapeGrid(g), g) for g in vs]
    grids = [(v, g) for v, g in grids if v is not None]
    if len(grids) < 2:
        return None
    G = np.stack([v for v, _ in grids])
    D = np.abs(G[:, None, :]-G[None, :, :]).sum(-1)
    med = grids[int(np.argmin(D.sum(1)))][1]
    return dict(strokes=[ST.SmoothStroke(s) for s in med["strokes"]],
                priorD=round(float(med.get("priorD", 0.5)), 4),
                advance=round(float(med.get("advance", 0.9)), 3),
                lead=round(float(med.get("lead", 0.0)), 3),
                width=round(float(med.get("width", 0.9)), 3),
                top=round(float(med.get("top", 1.0)), 3),
                bot=round(float(med.get("bot", 0.0)), 3),
                entryY=round(float(med.get("entryY", 0.3)), 3),
                exitY=round(float(med.get("exitY", 0.3)), 3),
                nPool=len(vs))
\end{lstlisting}
\Q{Justification}{Interpret the method rather than infer a benchmark}
A medoid minimises summed within-pool L1 descriptor distance and remains an actual sample. A denoised prototype combines compatible resampled strokes; it is a separate constructed object. \cite{code2026,haines2016}.


\section{Prototype}\label{fn:curate:Prototype}
\begin{lstlisting}[style=sig]
def Prototype(ch, vs)
\end{lstlisting}
\Q{Evidence}{Source and method link}
\path{Software/DemoCode/curate.py} line 173. Implementation evidence: \cite{code2026}. Related method, alternatives and library discussion: Section~\ref{method:medoid}.


\Q{Purpose}{What this contributes}
Choose representative writer glyphs while rejecting fragmented outliers.


\Q{Method}{Current implementation}
Choose representative writer glyphs while rejecting fragmented outliers. The excerpt below is from the current definition; the complete function is in the source snapshot. Its constants and branches, rather than an older signature, govern this path.


\begin{lstlisting}[style=plain]
def Prototype(ch, vs):
    proto = ST._DenoisedPrototype(vs)
    if proto is None:
        return None
    pts = [p for s in proto for p in s]
    xs = [p[0] for p in pts]
    ys = [p[1] for p in pts]
    return dict(strokes=[ST.SmoothStroke(s) for s in proto],
                priorD=round(float(np.median([g.get("priorD", 0.5) for g in vs])), 4),
                advance=round(float(np.median([g["advance"] for g in vs])), 3),
                lead=round(float(np.median([g.get("lead", 0.0) for g in vs])), 3),
                width=round(max(xs)-min(xs), 3),
                top=round(max(ys), 3), bot=round(min(ys), 3),
                entryY=round(float(np.median([g.get("entryY", 0.3) for g in vs])), 3),
                exitY=round(float(np.median([g.get("exitY", 0.3) for g in vs])), 3),
                nPool=len(vs))
\end{lstlisting}
\Q{Justification}{Interpret the method rather than infer a benchmark}
A medoid minimises summed within-pool L1 descriptor distance and remains an actual sample. A denoised prototype combines compatible resampled strokes; it is a separate constructed object. \cite{code2026,haines2016}.


\section*{Stage: Trajectory rasterisation}
\section{Method background: Trajectory rasterisation}\label{method:render}
\Q{Purpose}{The role in this system}
Make an image of a generated path for preview and recognition-based scoring.


\Q{Justification}{The reason for this construction}
Sampling straight segments and applying a controlled width makes candidate rendering explicit. The uniform width models one drawing pen rather than source ink intensity.


\Q{Alternatives}{What changes if another method is used}
Vector display keeps exact paths for viewing but classifiers need rasters. Antialiased line drawing and supersampling change edge pixels and cost. Rendering width is not physical pressure.


\Q{Library}{Implementation speed and accuracy}
Pillow ImageDraw or OpenCV polylines can rasterise paths. Coverage rules, antialiasing, round caps and blur alter classifier inputs. Pixel similarity and performance should be compared separately. \cite{pillowdraw,opencvdraw,code2026}.


\section{DrawGlyph}\label{fn:curate:DrawGlyph}
\begin{lstlisting}[style=sig]
def DrawGlyph(g, cell=150, pad=14)
\end{lstlisting}
\Q{Evidence}{Source and method link}
\path{Software/DemoCode/curate.py} line 67. Implementation evidence: \cite{code2026}. Related method, alternatives and library discussion: Section~\ref{method:render}.


\Q{Purpose}{What this contributes}
Make an image of a generated path for preview and recognition-based scoring.


\Q{Method}{Current implementation}
Make an image of a generated path for preview and recognition-based scoring. The excerpt below is from the current definition; the complete function is in the source snapshot. Its constants and branches, rather than an older signature, govern this path.


\begin{lstlisting}[style=plain]
def DrawGlyph(g, cell=150, pad=14):
    from PIL import Image, ImageDraw
    im = Image.new("L", (cell, cell), 255)
    d = ImageDraw.Draw(im)
    strokes = g.get("strokes") or []
    pts = [p for s in strokes for p in s]
    if not pts:
        return im
    xs = [p[0] for p in pts]
    ys = [p[1] for p in pts]
    w = max(1e-6, max(xs)-min(xs))
    h = max(1e-6, max(ys)-min(ys))
    k = (cell-2*pad)/max(w, h)
    ox = (cell-w*k)/2.0
    oy = (cell-h*k)/2.0
    for s in strokes:
        q = [(ox+(p[0]-min(xs))*k, cell-(oy+(p[1]-min(ys))*k)) for p in s]
        if len(q) > 1:
            d.line(q, fill=0, width=2)
    return im
\end{lstlisting}
\Q{Justification}{Interpret the method rather than infer a benchmark}
Sampling straight segments and applying a controlled width makes candidate rendering explicit. The uniform width models one drawing pen rather than source ink intensity. \cite{pillowdraw,opencvdraw,code2026}.


\section{Sheet}\label{fn:curate:Sheet}
\begin{lstlisting}[style=sig]
def Sheet(writer, curated=False, cell=150)
\end{lstlisting}
\Q{Evidence}{Source and method link}
\path{Software/DemoCode/curate.py} line 95. Implementation evidence: \cite{code2026}. Related method, alternatives and library discussion: Section~\ref{method:render}.


\Q{Purpose}{What this contributes}
One image per character, every variant numbered the number under a variant is its index, which is what drop and keep take.  a variant marked proto is the denoised prototype the builder prepended rather than one the writer actually wrote.


\Q{Method}{Current implementation}
Make an image of a generated path for preview and recognition-based scoring. The excerpt below is from the current definition; the complete function is in the source snapshot. Its constants and branches, rather than an older signature, govern this path.


\begin{lstlisting}[style=plain]
def Sheet(writer, curated=False, cell=150):
    from PIL import Image, ImageDraw
    prof, path = LoadProfile(writer, curated)
    glyphs = prof.get("glyphs") or {}
    outDir = OUT/writer
    outDir.mkdir(parents=True, exist_ok=True)
    made = []
    for ch in sorted(glyphs):
        vs = glyphs[ch]
        cols = min(10, max(1, len(vs)))
        rows = (len(vs)+cols-1)//cols
        im = Image.new("L", (cols*cell, rows*(cell+22)+26), 255)
        d = ImageDraw.Draw(im)
        d.text((6, 6), "%s  %r  %d variants" % (writer, ch, len(vs)), fill=0)
        for i, g in enumerate(vs):
            cx, cy = (i % cols)*cell, 26+(i//cols)*(cell+22)
            im.paste(DrawGlyph(g, cell), (cx, cy))
            tag = "%d%s" % (i, "  proto" if g.get("denoised") else "")
            d.text((cx+4, cy+cell+4), tag, fill=0)
            d.text((cx+4, cy+cell+12), "pd %.2f" % float(g.get("priorD", 0.5)), fill=0)
        safe = "".join(c if c.isalnum() else "_%02x" % ord(c) for c in ch)
        f = outDir/("%s.png" % safe)
        im.save(f)
        made.append((ch, len(vs), f.name))
    index = outDir/"INDEX.txt"
    with open(index, "w", encoding="utf-8") as fh:
        fh.write("%s, from %s\n\n" % (writer, path.name))
        fh.write("the number under a variant is its index, for --drop and --keep\n")
        fh.write("pd is the prior distance: lower means closer to the usual shape\n\n")
        for ch, n, name in made:
            fh.write("%-4r %3d variants  %s\n" % (ch, n, name))
    print("wrote %d character sheets to %s" % (len(made), outDir))
    print("  index: %s" % index)
    return made
\end{lstlisting}
\Q{Justification}{Interpret the method rather than infer a benchmark}
Sampling straight segments and applying a controlled width makes candidate rendering explicit. The uniform width models one drawing pen rather than source ink intensity. \cite{pillowdraw,opencvdraw,code2026}.

